\documentclass[11pt,a4paper]{article}

\usepackage[utf8]{inputenc}
\usepackage[T1]{fontenc}
\usepackage[polish]{babel}
\usepackage[margin=2.5cm]{geometry}
\usepackage{amsmath,amssymb}
\usepackage{graphicx}
\usepackage{booktabs}
\usepackage{longtable}
\usepackage{array}
\usepackage{xcolor}
\usepackage{placeins}
\usepackage{hyperref}

\hypersetup{colorlinks=true, linkcolor=blue, urlcolor=blue, citecolor=blue}

\title{%
    Przygotowanie benchmarku i protokołu eksperymentalnego do oceny wpływu\\
    topologii sprzętu oraz kodowania logicznego na jakość obliczeń kwantowych\\
    \large Sprawozdanie z Miesiąca 2
}
\author{Karina Leśkiewicz}
\date{Sierpień 2026}

\begin{document}

\maketitle
\thispagestyle{empty}

% =====================================================================
% UWAGA ROBOCZA (usunąć przed finalnym wysłaniem):
% Lista otwartych "kwiatków"/luk merytorycznych jest w project_review_gaps.md
% i MUSI zostać w całości zweryfikowana i domknięta przed wysłaniem pracy.
% =====================================================================
\begin{center}
\fbox{\parbox{0.92\textwidth}{\color{red}\small\textbf{NOTATKA ROBOCZA (usunąć przed wysłaniem):}
przed złożeniem pracy należy w~całości sprawdzić i~domknąć listę otwartych luk
merytorycznych w~pliku \texttt{project\_review\_gaps.md} (m.in. crossover metod,
koszt cross-topology Emeralda, korelacja SWAP$\leftrightarrow$degradacja, pasma driftu,
oraz niewykorzystane elementy gałęzi logicznej \texttt{[[4,2,2]]}).}}
\end{center}

% =====================================================================
\section{Wprowadzenie: druga metoda mitygacji --- Clifford Data Regression}
\label{sec:cdr-intro}
% =====================================================================

Miesiąc~2 rozszerza badanie z pojedynczej metody mitygacji do \emph{porównania} strategii
mitygacji i bliskoterminowej ochrony logicznej wzdłuż tej samej osi kosztu routingu. W~Miesiącu~1
potok stanowiła para REM${+}$ZNE; teraz dokładam drugą metodę mitygacji błędów ---
\textbf{Clifford Data Regression (CDR)} --- oraz warstwę „logical-ish'' opartą na detekcji błędów
(kod $[[4,2,2]]$, sekcja osobna). Wszystkie metody oceniam na tych samych metrykach jakości
(błąd $\langle O\rangle$, dokładność, mirror fidelity) oraz na \emph{wspólnej walucie kosztu} ---
liczbie shotów QPU potrzebnych do osiągnięcia zadanej tolerancji --- tak, aby były wprost
porównywalne mimo różnych narzutów natywnych.

\paragraph{Pozycjonowanie i wkład.} Samo \emph{istnienie} progu, przy którym mitygacja przestaje
się opłacać, nie jest wynikiem nowym --- wynika wprost z ram bias/wariancja~\cite{cai2023}
(rosnący narzut próbkowania $C_{em}$ w~ZNE) oraz z teorii jałowych płaskowyżów indukowanych
szumem~\cite{wang2021}: przy
dostatecznej głębokości stan zbiega do maksymalnie zmieszanego i \emph{każda} metoda musi zawieść.
Wkładem tej pracy nie jest więc odkrycie, że próg istnieje, lecz jego \textbf{ilościowy,
rozdzielony po urządzeniach i topologiach pomiar}: gdzie dokładnie leży punkt opłacalności --- na
gwieździe IQM~Spark (ODRA) wobec kraty kwadratowej (Garnet, Emerald) --- jest wielkością
specyficzną dla urządzenia, kalibracji i topologii, której nie sposób podać bez pomiaru.
Kluczowe przeramowanie polega na tym, że \textbf{osią kosztu jest narzut routingu (bramki SWAP),
a nie sama głębokość obwodu}: ten sam obwód logiczny przekracza próg w \emph{różnych} punktach na
różnych topologiach wyłącznie z powodu rozmieszczenia (placementu) kubitów. Widać to już w~samej
drabinie kosztu: identyczny obwód wzorca \texttt{star} kompiluje się do $12$~CZ przy rozmieszczeniu
zwartym, ale do $78$~CZ przy rozproszonym na Emeraldzie (sekcja mapowania). To porównanie --- próg
opłacalności mitygacji jako funkcja \emph{topologii i metody}, mierzony wzdłuż osi narzutu
routingu --- stanowi właściwy wkład sprawozdania; wynik progowy z Miesiąca~1 (${\approx}85$--$90$~CZ
na gwieździe) jest jego punktem odniesienia, a nie tezą samą w sobie.

\paragraph{Motywacja: inny profil narzutu niż ZNE.} ZNE sztucznie zwiększa szum i ekstrapoluje do
zera, więc jego estymator jest różnicą dużych liczb o naprzemiennych znakach; narzut wariancji
$C_{em}$ rośnie z kosztem obwodu i to on wyznacza górną granicę użyteczności ZNE (w~Miesiącu~1
zmierzoną w okolicach $85$--$90$~CZ, gdzie współczynnik redukcji błędu spadał do ${\approx}1{,}0$).
CDR \emph{nie} zwiększa szumu --- zamiast tego \emph{uczy się} korekty z danych --- więc ma z zasady
inny profil załamania. Jest to naturalny kandydat na drugą metodę: pytanie badawcze brzmi, która
z metod załamuje się pierwsza wraz z rosnącym kosztem routingu.

\paragraph{Mechanizm.} CDR~\cite{czarnik2021} (Czarnik i~in., \emph{Error mitigation with Clifford
quantum-circuit data}, Quantum~\textbf{5}, 592, 2021; arXiv:2005.10189) buduje klasyczny model korekcji na
\textbf{obwodach treningowych bliskich Cliffordowi}. Konstrukcja przebiega w trzech krokach:
\begin{enumerate}
\item \textbf{Generacja zbioru treningowego.} Z obwodu docelowego tworzę warianty o
      \emph{identycznej strukturze} (ten sam wzorzec splątający, ta sama transpilacja i routing),
      w których obroty niekliffordowskie --- w naszym VQC data-zależne $R_y(x)$ --- zostają
      „przyciągnięte'' do najbliższego kąta kliffordowskiego (wielokrotności $\pi/2$). Dzięki
      twierdzeniu Gottesmana--Knilla wartość oczekiwana takiego obwodu jest \emph{dokładnie i
      efektywnie} obliczalna symulacją stabilizatorową, co daje bezszumową referencję
      $\langle O\rangle_{\text{ideal}}$ bez kosztu wektora stanu.
\item \textbf{Zbieranie par (szum, ideał).} Każdy obwód treningowy wykonuję na urządzeniu (lub w
      kalibrowanym symulatorze), otrzymując parę
      $\bigl(\langle O\rangle_{\text{noisy}},\,\langle O\rangle_{\text{ideal}}\bigr)$. Ponieważ
      obwody treningowe mają ten sam profil szumu co docelowy, para ta próbkuje odwzorowanie
      szumu urządzenia.
\item \textbf{Dopasowanie i zastosowanie.} Dopasowuję regresję (w wariancie podstawowym
      \emph{liniową})
      \begin{equation}
      \langle O\rangle_{\text{ideal}} \approx a\,\langle O\rangle_{\text{noisy}} + b,
      \end{equation}
      a wyuczone współczynniki $(a,b)$ stosuję do surowej wartości obwodu \emph{docelowego}. Estymator
      CDR nie wymaga wzmacniania szumu, więc nie dziedziczy eksplozji wariancji ZNE.
\end{enumerate}

\paragraph{Koszt i założenia.} Narzutem CDR jest budżet obwodów treningowych: $N$ obwodów bliskich
Cliffordowi ${\times}$ liczba shotów każdego. Ten koszt wliczam do wspólnej waluty
(shots-to-target), aby porównanie z ZNE było uczciwe. Kluczowe założenie CDR to
\emph{stacjonarność i wyuczalność} odwzorowania szumu: jeśli kalibracja dryfuje między pomiarem
zbioru treningowego a obwodu docelowego, wyuczona korekta przestaje pasować. Dlatego --- inaczej
niż w ZNE --- obwody treningowe i docelowy uruchamiam w tym samym oknie/zadaniu, a wrażliwość CDR na
dryf (udokumentowany w sekcji~\ref{sec:cv-full}) traktuję jako część wyniku, a nie usterkę
pomiaru.

\FloatBarrier

% =====================================================================
\section{Pełna walidacja krzyżowa 5-fold na QPU (Banknote, wszystkie głębokości)}
\label{sec:cv-full}
% =====================================================================

Rozszerzyłam walidację krzyżową z Miesiąca~1 (dotąd tylko $L=6$, dwa ansätze) na
\textbf{pełną siatkę}: trzy ansätze $\times$ trzy głębokości $\{2,4,6\}$ $\times$ wszystkie
pięć foldów, wykonane na rzeczywistym QPU ODRA (IQM~Spark). Metodyka jest identyczna z pracą
źródłową i z Miesiącem~1: dla każdego foldu używam \emph{jego własnych} wag (uczonych na
zbiorze treningowym tego foldu) i oceniam na jego rozłącznym zbiorze testowym --- pełnym
foldzie (274--275~pkt, naturalna proporcja klas), więc bez wycieku danych i z sensowną
macro-F1. Konfiguracja: DD wyłączone, 4096~shotów, potok raw${}\to{}$+REM${}\to{}$+ZNE
(Richardson, skale $\{1,3,5\}$). Wynik podaję jako średnią${}\pm{}$odchylenie
\emph{międzyfoldowe} (ddof${}=1$). Dane: \texttt{banknote\_hardware\_cv\_full.csv};
agregacja: \texttt{run\_cv\_banknote.py -\/-aggregate-csv}.

\begin{table}[htbp]
\centering
\small
\caption{Banknote na ODRA --- \textbf{pełna 5-fold CV}, wszystkie ansätze i głębokości
(DD-off, 4096 shots, pełny fold). Dokładność i macro-F1 podane jako
średnia${}\pm{}$odchylenie międzyfoldowe; kolumna err to średni $|\langle O\rangle$-błąd$|$
wobec ideału bezszumowego. ZNE ${}={}$ REM${+}$ekstrapolacja Richardsona. Wiersze uporządkowane
rosnącym kosztem CZ w obrębie każdej głębokości (drabina routingu star${}<{}$odra${}<{}$sim).}
\label{tab:cv-full}
\begin{tabular}{@{}lc cc cc c@{}}
\toprule
& & \multicolumn{2}{c}{\textbf{dokładność}} & \multicolumn{2}{c}{\textbf{macro-F1}}
  & err $\langle O\rangle$ \\
\cmidrule(lr){3-4}\cmidrule(lr){5-6}
Ansatz & CZ & raw & ZNE & raw & ZNE & raw${}\to{}$ZNE \\
\midrule
\multicolumn{7}{@{}l}{\emph{głębokość 2}} \\
star      & 8  & $0{,}840\pm0{,}035$ & $0{,}843\pm0{,}031$ & $0{,}834\pm0{,}041$ & $0{,}838\pm0{,}033$ & $0{,}072\to0{,}043$ \\
odra      & 13 & $0{,}855\pm0{,}029$ & $0{,}875\pm0{,}020$ & $0{,}854\pm0{,}029$ & $0{,}873\pm0{,}021$ & $0{,}112\to0{,}058$ \\
simulator & 20 & $0{,}808\pm0{,}038$ & $0{,}813\pm0{,}079$ & $0{,}796\pm0{,}050$ & $0{,}792\pm0{,}098$ & $0{,}149\to0{,}163$ \\
\midrule
\multicolumn{7}{@{}l}{\emph{głębokość 4}} \\
star      & 16 & $0{,}886\pm0{,}009$ & $0{,}886\pm0{,}035$ & $0{,}885\pm0{,}010$ & $0{,}881\pm0{,}039$ & $0{,}126\to0{,}036$ \\
odra      & 36 & $0{,}682\pm0{,}187$ & $0{,}744\pm0{,}100$ & $0{,}645\pm0{,}223$ & $0{,}729\pm0{,}107$ & $0{,}259\to0{,}217$ \\
simulator & 52 & $0{,}699\pm0{,}074$ & $0{,}711\pm0{,}052$ & $0{,}672\pm0{,}089$ & $0{,}677\pm0{,}089$ & $0{,}246\to0{,}220$ \\
\midrule
\multicolumn{7}{@{}l}{\emph{głębokość 6}} \\
star      & 24 & $0{,}868\pm0{,}058$ & $0{,}886\pm0{,}039$ & $0{,}862\pm0{,}071$ & $0{,}883\pm0{,}044$ & $0{,}161\to0{,}093$ \\
odra      & 58 & $0{,}731\pm0{,}050$ & $0{,}776\pm0{,}063$ & $0{,}698\pm0{,}062$ & $0{,}760\pm0{,}068$ & $0{,}226\to0{,}177$ \\
simulator & 86 & $0{,}583\pm0{,}062$ & $0{,}594\pm0{,}075$ & $0{,}537\pm0{,}084$ & $0{,}548\pm0{,}098$ & $0{,}285\to0{,}281$ \\
\bottomrule
\end{tabular}
\end{table}

\FloatBarrier

% ---------------------------------------------------------------------
\subsection{Wnioski}
% ---------------------------------------------------------------------

\begin{enumerate}

\item \textbf{Drabina routingu obowiązuje na każdej głębokości.} Dokładność surowa spada
      monotonicznie z kosztem CZ przy \emph{każdej} głębokości i zachowuje ten sam porządek
      ansätze (dla $L=6$: star~$0{,}868$ ${>}$ odra~$0{,}731$ ${>}$ simulator~$0{,}583$). Jest
      to najczystsze jak dotąd potwierdzenie, że transpilowany koszt CZ jest właściwą zmienną
      niezależną --- teraz z pasmami błędu z pięciu foldów.

\item \textbf{ZNE odzyskuje dokładność dokładnie tam, gdzie istnieje luka --- i wygasa na
      szczycie drabiny.} Największy zysk mają średnio- i wysokokosztowe pierścienie
      (odra $L{=}4$: $0{,}682\to0{,}744$; odra $L{=}6$: $0{,}731\to0{,}776$; ${+}0{,}05$--$0{,}06$).
      Dla tanich obwodów dokładność surowa jest już przy ideale, więc ZNE nic nie wnosi
      (star $L{=}2,4$). W najdroższym punkcie --- \textbf{simulator $L{=}6$, 86~CZ} --- odzysk
      \emph{zanika}: $0{,}583\to0{,}594$ (w granicach rozrzutu), a błąd $\langle O\rangle$
      $0{,}285\to0{,}281$ (współczynnik ${\approx}1{,}0$). To bezpośredni pomiar granicy
      opłacalności ZNE w okolicach $85$--$90$~CZ zapowiadanej w Miesiącu~1.

\item \textbf{Redukcja błędu $\langle O\rangle$ jest monotoniczna w koszcie i najsilniejsza przy
      tanich obwodach} (star $L{=}4$: $0{,}126\to0{,}036$, ${\approx}3{,}5\times$; odra $L{=}2$:
      ${\approx}1{,}9\times$; simulator $L{=}6$: ${\approx}1{,}0\times$, martwa). Potwierdza to,
      że siła ZNE słabnie z routingiem niezależnie od tego, czy przekłada się na dokładność.

\item \textbf{Nadal nie odtwarzamy zapaści pierścieni z pracy źródłowej.} Najgorszy przypadek
      (simulator $L{=}6$) degraduje się w stronę floor klasy większościowej ($0{,}56$),
      osiągając F1~$0{,}537$, ale \emph{nie} zapada do F1${\approx}0{,}04$ raportowanego w pracy;
      odra $L{=}6$ pozostaje zdrowy (F1~$0{,}698$). W żadnym wierszu DD-off macro-F1 nie odrywa
      się od dokładności --- brak zapaści klasowej.

\item \textbf{Pierścień odra przy $L{=}4$ jest niestabilny międzyfoldowo.} Dokładność
      $0{,}682\pm0{,}187$ i F1 $0{,}645\pm0{,}223$ mają ogromny rozrzut: część foldów utrzymuje
      wynik, część się załamuje. Sam ten rozrzut jest wynikiem --- wskazuje na indukowaną
      routingiem niestabilność, a nie stabilny punkt pracy, i jest kandydatem do dokładniejszego
      zbadania (drift, dobór kubitów) na Garnecie/Emeraldzie.

\item \textbf{Dryf urządzenia względem Miesiąca~1.} Te same wagi, ansätze i transpilacja, lecz
      pomiar $L{=}6$ z pełnej CV (12.08) daje niższe wartości niż CV z Miesiąca~1 (03.08):
      odra $0{,}809\pm0{,}042 \to 0{,}731\pm0{,}050$, simulator $0{,}626\pm0{,}091 \to
      0{,}583\pm0{,}062$. Różnicę --- przy identycznym potoku --- przypisuję dryfowi kalibracji
      między dniami; jest to bezpośrednie, ilościowe potwierdzenie tezy o zależności zachowania
      on-device od stanu urządzenia. Wynikowo trend (drabina, okno ZNE, brak zapaści klasowej)
      jest niezmieniony; przesuwa się jedynie bezwzględny poziom.

\end{enumerate}

% =====================================================================
\section{Porównanie metod na symulatorze ODRA (\texttt{two\_curves})}
\label{sec:method-comparison}
% =====================================================================

Poniżej zestawiam trzy metody obsługi błędów --- ZNE, CDR i detekcję $[[4,2,2]]$ --- wzdłuż tej
samej osi kosztu routingu, na wiernym symulatorze gwiazdy ODRA z \emph{realną} kalibracją
(\texttt{calibrations/odra.json}). Ustawienia dobrane pod wiarygodność: 8192 shotów (błąd
statystyczny $\langle Z\rangle$ ${\approx}0{,}005$), 5 powtórzeń dla pasm błędu (ZNE), 200 obwodów
treningowych (CDR), pełny 30-punktowy zbiór ewaluacyjny, $p_{2q}=0{,}019$ (zmierzony błąd 2q
ODRA) dla $[[4,2,2]]$. Emeralda ($2^{54}$ stanów) \emph{nie da się} zasymulować klasycznie ---
cross-topology musi przyjść ze sprzętu; symulator ODRA jest wierną kontrolą.

\begin{table}[htbp]
\centering
\small
\caption{ZNE: ablacja raw${}\to{}$DD${+}$REM${}\to{}$ZNE (Richardson), pasma $\pm$odchylenie po
5 powtórzeniach. Accuracy i macro-F1 (zbiór zbalansowany, więc F1${}\approx{}$acc).}
\begin{tabular}{@{}lc ccc cc@{}}
\toprule
Wzorzec & CZ & raw & DD${+}$REM & ZNE & acc (raw${\to}$ZNE) & F1 (raw${\to}$ZNE) \\
\midrule
star        & 12 & $0{,}193$ & $0{,}162$ & $\mathbf{0{,}024}$ & $0{,}87{\to}0{,}87$ & $0{,}87{\to}0{,}87$ \\
linear      & 27 & $0{,}408$ & $0{,}378$ & $0{,}243$ & $0{,}89{\to}0{,}90$ & $0{,}89{\to}0{,}90$ \\
circular    & 36 & $0{,}405$ & $0{,}382$ & $0{,}305$ & $0{,}90{\to}0{,}90$ & $0{,}90{\to}0{,}90$ \\
all\_to\_all & 54 & $0{,}454$ & $0{,}433$ & $0{,}345$ & $0{,}90{\to}0{,}87$ & $0{,}90{\to}0{,}87$ \\
\bottomrule
\end{tabular}
\end{table}

\begin{table}[htbp]
\centering
\small
\caption{CDR: błąd $\langle O\rangle$ raw${}\to{}$CDR, jakość dopasowania $r^2$, accuracy i macro-F1.}
\begin{tabular}{@{}lc cc c cc@{}}
\toprule
Wzorzec & CZ & raw & CDR & $r^2$ & acc (raw${\to}$CDR) & F1 (raw${\to}$CDR) \\
\midrule
star        & 12 & $0{,}172$ & $\mathbf{0{,}028}$ & $0{,}98$ & $0{,}87{\to}0{,}87$ & $0{,}87{\to}0{,}87$ \\
linear      & 27 & $0{,}404$ & $0{,}228$ & $0{,}79$ & $0{,}90{\to}0{,}90$ & $0{,}90{\to}0{,}90$ \\
circular    & 36 & $0{,}398$ & $0{,}199$ & $0{,}66$ & $0{,}90{\to}0{,}90$ & $0{,}90{\to}0{,}90$ \\
all\_to\_all & 54 & $0{,}448$ & $0{,}278$ & $0{,}49$ & $0{,}90{\to}0{,}90$ & $0{,}90{\to}0{,}90$ \\
\bottomrule
\end{tabular}
\end{table}

\begin{table}[htbp]
\centering
\small
\caption{$[[4,2,2]]$ detekcja (Option C, mirror): logiczny błąd $\langle O\rangle$ (bez${}\to{}$z
detekcją; ideał $\bar Z_1=+1$), fidelity oraz współczynnik akceptacji (koszt post-selekcji).}
\begin{tabular}{@{}cc cc cc c@{}}
\toprule
Warstwy & CZ & \multicolumn{2}{c}{logiczny $\langle O\rangle$-błąd} &
  \multicolumn{2}{c}{fidelity} & akceptacja \\
\cmidrule(lr){3-4}\cmidrule(lr){5-6}
 & & bez & z detekcją & bez & z detekcją & \\
\midrule
1 & 34  & $0{,}597$ & $\mathbf{0{,}419}$ & $0{,}396$ & $0{,}637$ & $0{,}62$ \\
2 & 63  & $0{,}874$ & $0{,}788$ & $0{,}226$ & $0{,}410$ & $0{,}55$ \\
3 & 104 & $0{,}978$ & $0{,}937$ & $0{,}151$ & $0{,}296$ & $0{,}51$ \\
4 & 157 & $0{,}987$ & $0{,}987$ & $0{,}131$ & $0{,}261$ & $0{,}51$ \\
\bottomrule
\end{tabular}
\end{table}

\begin{figure}[htbp]
\centering
\includegraphics[width=0.98\textwidth]{figures/fig_method_comparison.png}
\caption{Porównanie metod na osi kosztu routingu (ODRA sim, realna kalibracja).
(A) błąd $\langle O\rangle$ vs koszt CZ: raw (kreskowany) i mitygowany (ciągły) dla ZNE, CDR
oraz logiczny błąd $[[4,2,2]]$; wszystkie metody mitygują najmocniej przy tanim routingu i
słabną z kosztem. (B) efektywne shoty (skala log) na najtańszym szczeblu --- CDR osiąga ten sam
błąd co ZNE, lecz ${\sim}67{\times}$ drożej.}
\label{fig:method-comparison}
\end{figure}

\FloatBarrier
\paragraph{Wnioski.}
\begin{enumerate}
\item \textbf{Wszystkie trzy metody mitygują najmocniej przy tanim routingu i słabną
      przewidywalnie z kosztem CZ.} ZNE i CDR tną błąd $\langle O\rangle$ dramatycznie przy
      gwieździe (ZNE $0{,}193{\to}0{,}024$; CDR $0{,}172{\to}0{,}028$), a zysk maleje monotonicznie
      do ${\sim}54$~CZ. Detekcja $[[4,2,2]]$ powtarza ten wzorzec w logicznym $\langle O\rangle$-błędzie:
      $0{,}597{\to}0{,}419$ ($-30\%$) przy 34~CZ, aż do zaniku przy 157~CZ (obwód w pełni
      zdepolaryzowany, $0{,}987{\to}0{,}987$).
\item \textbf{Accuracy i macro-F1 są nasycone.} Na zbalansowanym \texttt{two\_curves} obie
      trzymają się $0{,}87$--$0{,}90$ niezależnie od metody i szczebla, mimo redukcji błędu
      $\langle O\rangle$ nawet 6-krotnej --- bo predykcja to $\mathrm{sgn}\langle O\rangle$, a
      punkty leżą daleko od granicy decyzyjnej. \emph{Czułym wskaźnikiem jest ciągły
      $\langle O\rangle$-błąd, nie accuracy} (potwierdzenie tezy z Miesiąca~1). Przy all\_to\_all
      ZNE wręcz \emph{obniża} accuracy ($0{,}90{\to}0{,}87$) --- dodana wariancja przerzuca
      pojedynczy znak blisko sufitu.
\item \textbf{Na wspólnej walucie kosztu ZNE dominuje.} ZNE i CDR docierają do niemal tego samego
      błędu (a CDR bywa lekko lepszy przy dużym koszcie, bo nie wzmacnia szumu), lecz efektywny
      koszt to ZNE ${=}8192{\times}3={}$24\,576 shotów wobec CDR ${=}8192{\times}201={}$1\,646\,592
      --- \textbf{${\sim}67{\times}$ drożej za ten sam wynik} (rys.~\ref{fig:method-comparison}B).
      $r^2$ CDR spada $0{,}98{\to}0{,}49$ z kosztem --- własny sygnał załamania metody.
\item \textbf{$[[4,2,2]]$ kupuje wierność kosztem shotów.} Detekcja podnosi fidelity ${\sim}1{,}6{\times}$
      i tnie logiczny $\langle O\rangle$-błąd przy niskim koszcie, przy akceptacji ${\sim}0{,}5$
      (${\sim}2{\times}$ więcej shotów). Uwaga metodologiczna: accuracy jest \emph{nieinformatywna}
      dla kodu detekcyjnego na surogacie Clifforda (wartości logiczne kolapsują do $\pm1$ lub $0$),
      więc właściwymi metrykami są logiczny $\langle O\rangle$-błąd, fidelity i koszt akceptacji.
\end{enumerate}

% =====================================================================
\section{Circuit twirling: tłumienie błędów koherentnych (wyprowadzenie)}
% =====================================================================

\paragraph{Motywacja.} Błąd bramki opisuje kanał kwantowy (odwzorowanie CPTP)
$\mathcal{E}$. Rozróżniamy dwa reżimy. Błąd \emph{koherentny} jest sam unitarny,
$\mathcal{E}(\rho)=V\rho V^\dagger$ (np.\ systematyczna nad-rotacja
$V=e^{-i\theta Z/2}$): jest odwracalny, dodaje się w~\emph{amplitudzie}
(po $N$ bramkach prawdopodobieństwo błędu $\sim(N\theta)^2$) i~\emph{systematycznie
przesuwa} $\langle O\rangle$. Błąd \emph{stochastyczny} to mieszanina Pauliego
$\mathcal{E}(\rho)=\sum_k p_k P_k\rho P_k$: dodaje się w~\emph{prawdopodobieństwie}
($\sim Np$). Kluczowe: ZNE zakłada gładkie skalowanie szumu, a~CDR liniową relację
$\langle O\rangle_{\text{ideal}}\!\approx\!a\langle O\rangle_{\text{noisy}}+b$ ---
błędy koherentne łamią oba założenia.

\paragraph{Macierz procesu $\chi$.} Dowolny kanał zapisujemy w~bazie Pauliego~\cite{nielsen2010}:
\begin{equation}
\mathcal{E}(\rho)=\sum_{a,b}\chi_{ab}\,P_a\,\rho\,P_b .
\end{equation}
Dla kanału unitarnego z~$V=\sum_a c_a P_a$ mamy $\chi_{ab}=c_a c_b^{*}$ (iloczyn
zewnętrzny wektora współczynników). Diagonala $\chi_{aa}=|c_a|^2$ to stochastyczne
prawdopodobieństwa błędów Pauliego; elementy \emph{poza-diagonalne} $\chi_{ab}$
($a\neq b$) kodują część \emph{koherentną}. Przykład: $V=e^{-i\theta Z/2}
=\cos\tfrac\theta2\,I-i\sin\tfrac\theta2\,Z$ daje
$\chi_{II}=\cos^2\tfrac\theta2$, $\chi_{ZZ}=\sin^2\tfrac\theta2$ oraz
niezerowe $\chi_{IZ}=\chi_{ZI}^{*}=i\cos\tfrac\theta2\sin\tfrac\theta2$.

\paragraph{Twirl Pauliego zeruje poza-diagonalę.} Twirlowany kanał to średnia po
grupie Pauliego~\cite{emerson2007,wallman2016}:
\begin{equation}
\mathcal{E}_T(\rho)=\frac{1}{4^n}\sum_k P_k^\dagger\,\mathcal{E}(P_k\rho P_k)\,P_k .
\end{equation}
Korzystając z~$P_kP_aP_k=(-1)^{\langle a,k\rangle}P_a$ (gdzie $\langle a,k\rangle\in\{0,1\}$
to symplektyczny iloczyn: $0$ dla komutacji, $1$ dla antykomutacji):
\begin{equation}
\mathcal{E}_T(\rho)=\sum_{a,b}\chi_{ab}
\underbrace{\Big[\tfrac{1}{4^n}\textstyle\sum_k(-1)^{\langle a\oplus b,\,k\rangle}\Big]}_{=\;\delta_{ab}}
P_a\rho P_b=\sum_a\chi_{aa}\,P_a\rho P_a .
\end{equation}
Suma wewnętrzna to ortogonalność charakterów: znika, chyba że $a=b$. Intuicyjnie ---
pole $\chi_{aa}$ mnoży się przez $(\pm1)^2\!=\!+1$ (zawsze przeżywa), pole
$\chi_{ab}$ ($a\neq b$) przez naprzemienne $\pm1$, które sumują się do zera. Zatem
twirl \textbf{kasuje część koherentną}, zostawiając czysto stochastyczny kanał
Pauliego. Nad-rotacja $\theta$ staje się losowym $Z$-flipem o~prawdopodobieństwie
$\sin^2\tfrac\theta2$, a~akumulacja błędu spada z~$\sim(N\theta)^2$ do~$\sim N\theta^2/4$.

\paragraph{Randomized compiling (implementacja).} Kanału nie uśrednia się wprost ---
\emph{randomizujemy obwód}~\cite{wallman2016}: każdą \emph{Cliffordowską} bramkę 2-kubitową
$G$ zastępujemy $B\,G\,A$ z~losową Paulią $A$ i~kompensacją $B=GAG^\dagger$; wtedy
$BGA=G$ (ideał zachowany), a~szum na $G$ jest twirlowany. $B$ jest Paulią
\emph{tylko dlatego, że $G$ jest Cliffordem}. Bramki jednokubitowe (rotacje $R_x$,
non-Clifford) \emph{wchłaniają} losowe Pauli do warstwy jednokubitowej --- bez
dodatkowej głębi --- i~\emph{nie} są twirlowane bezpośrednio (dla non-Clifforda
$GAG^\dagger$ nie jest Paulią). Uśrednienie $\langle O\rangle$ po $M$ losowych
instancjach realizuje kanał twirlowany. Reżim eksperymentalny (odniesienie): RC na
sprzęcie nadprzewodzącym stosuje $M\!\sim\!20$ randomizacji ($5$–$10$ wystarcza do
zbieżności) i $\sim\!8192$ shotów na obwód~\cite{hashim2021}; CDR typowo $\sim\!10$
obwodów treningowych near-Clifford~\cite{czarnik2021}. Nasza kampania celuje w ten
reżim w pełnej wersji (metodologia §9b); pilot używa świadomie okrojonych
$M\!=\!8$, shots $2048$.

\paragraph{Konsekwencja dla ZNE i CDR.} Po twirlingu szum jest Pauliego, więc
$\langle O\rangle(\lambda)$ maleje monotonicznie (poprawna ekstrapolacja ZNE), a~relacja
$\langle O\rangle_{\text{ideal}}$–$\langle O\rangle_{\text{noisy}}$ jest liniowa
(poprawny fit CDR). W~symulacji z~wstrzykniętą nad-rotacją $\theta=0{,}15$~rad/CX
twirl+ZNE poprawia $\langle O\rangle$-błąd $\sim2{,}2\times$ względem samego ZNE, a~twirl+CDR
przywraca $r^2$ dopasowania z~$0{,}81$ do~$0{,}99$ (błąd $0{,}109\!\to\!0{,}015$).
To sugeruje, że obserwowany \emph{spadek} $r^2$ CDR z~kosztem 2q na sprzęcie
(rozdz.~o~porównaniu metod) jest sygnaturą błędu koherentnego, którą twirling
powinien odwrócić --- hipotezę tę zweryfikowaliśmy bezpośrednio na Garnecie (poniżej).

\paragraph{Pilot na sprzęcie (Garnet).} IQM Resonance nie wystawia twirlingu serwerowo (tylko
dynamical decoupling), więc realizujemy go \emph{client-side}: losowe Pauli wokół każdej bramki
dwukubitowej, uśredniane po $M$ instancjach. Drabina trzech punktów kosztu routingu (12/57/96
transpiled CZ na modelu \texttt{two\_curves} offset $0{,}1$), $n_{\text{train}}=10$, $M=6$, $3$
punkty ewaluacji, $1024$ shoty. Dla CDR (najostrzejsza metryka --- jakość dopasowania $r^2$):
\begin{center}
\begin{tabular}{cccc}
\toprule
transpiled CZ & $r^2$ (bez twirl) & $r^2$ (twirl) & $\Delta r^2$ \\
\midrule
12 (optimal, star)       & 0{,}957 & 0{,}986 & $+0{,}029$ \\
57 (optimal, all-to-all) & 0{,}820 & 0{,}939 & $+0{,}119$ \\
96 (adverse, all-to-all) & 0{,}837 & 0{,}909 & $+0{,}072$ \\
\bottomrule
\end{tabular}
\end{center}
$r^2$ bez twirlingu spada z~kosztem routingu (sygnatura błędu koherentnego łamiącego liniowość
CDR), a~twirling przywraca je na~każdym punkcie, z~zyskiem rosnącym przy wysokim CZ --- co
potwierdza mechanizm na~prawdziwym sprzęcie; błąd $\langle O\rangle$ CDR maleje po~twirlingu na
każdym punkcie (np.\ $0{,}069\!\to\!0{,}056$ przy 96 CZ). Dla ZNE wynik jest \emph{mieszany}
(twirling pomaga przy 12 CZ: $0{,}083\!\to\!0{,}048$, ale szkodzi przy 57/96 CZ), co przypisujemy
osłabieniu twirlingu logicznego pod foldingiem ZNE oraz ograniczonej statystyce pilota --- wątek
do rozstrzygnięcia w~wersji pełnej (twirling natywny, unitarity RB, pasma driftu; metodologia §9).

\paragraph{Hipotezy.} Wyniki (sim + pilot sprzętowy) prowadzą do następujących hipotez, które
wersja pełna ma zweryfikować ilościowo:
\begin{itemize}
\item[\textbf{H1}] Narzut routingu (liczba CZ) akumuluje \emph{koherentny} błąd, który łamie
  liniowość dopasowania CDR --- $r^2$ CDR maleje z~kosztem 2q. \emph{(Wsparte: $r^2$ spada
  $0{,}98\!\to\!0{,}49$ z~kosztem; pilot $0{,}96\!\to\!0{,}82$--$0{,}84$.)}
\item[\textbf{H2}] Circuit twirling zamienia błąd koherentny w~stochastyczny i~\emph{przywraca}
  liniowość CDR ($r^2$ rośnie), z~zyskiem rosnącym z~kosztem CZ. \emph{(Potwierdzone na Garnecie:
  $\Delta r^2 = +0{,}03$/$+0{,}12$/$+0{,}07$ przy 12/57/96 CZ.)}
\item[\textbf{H3}] Zysk twirlingu jest proporcjonalny do \emph{frakcji koherentnej} szumu ---
  przewidywalny z~unitarity $u$ (mierzonej przez URB). \emph{(Do sprawdzenia: predicted$=$measured.)}
\item[\textbf{H4}] CDR jest bardziej wrażliwe na koherencję niż ZNE (dopasowanie liniowe pęka;
  ekstrapolacja ZNE jest odporniejsza) $\Rightarrow$ twirling pomaga CDR bardziej niż ZNE.
  \emph{(Wsparte: pilot CDR jednoznaczny, ZNE mieszany.)}
\item[\textbf{H5}] twirling$+$ZNE składa się, gdy frakcja koherentna jest wysoka (sim: 6/6 seedów
  poprawy); na sprzęcie opłaca się tam, gdzie koherencja dominuje.
\end{itemize}

\paragraph{Decyzja: kampania twirlingu skupia się na CDR (nie ZNE).} Podkreślamy, że
\emph{nie} rezygnujemy z~ZNE jako metody --- ZNE pozostaje na~osi porównania metod
(crossover ZNE/CDR). Rezygnujemy z~\emph{kombinacji} twirling$+$ZNE jako osobnego wyniku,
z~trzech spójnych powodów: (i)~ekstrapolacja ZNE jest z~natury \emph{odporna} na~błąd
koherentny (nie opiera się na~dopasowaniu liniowym, w~przeciwieństwie do~CDR, którego fit
liniowy koherencja łamie --- H4), więc twirling ma tam mało do~zyskania; (ii)~pomiar URB$+$RB
wskazuje, że frakcja koherentna \emph{na~poziomie bramki} jest mała i~utopiona w~drifcie
($u-p^2\approx0{,}02$ przy drifcie $\approx0{,}09$), a~sim pokazuje zysk twirling$+$ZNE tylko
przy \emph{dużej} koherencji --- czego Garnet nie ma; (iii)~na~sprzęcie twirling$+$ZNE
\emph{szkodził} przy wysokim koszcie 2q (narzut twirlingu $>$ zysk, gdy koherencji mało).
Wniosek: twirling wnosi wartość tam, gdzie metoda jest \emph{wrażliwa na~koherencję} (CDR), a~nie
tam, gdzie jest na~nią odporna (ZNE). Kampania pełna (§9b) mierzy więc twirling$+$CDR z~pasmami
driftu oraz koherencję \emph{na~poziomie obwodu} (właściwy predyktor), pomijając kosztowne
gate-URB (drift zaciera sygnał).

\paragraph{Plan weryfikacji i koszt (wersja pełna, §9b).} (1) \textbf{Unitarity RB} na krawędziach
Garneta (opt.\ i~adverse) $\to$ zmierzona frakcja koherentna $u$ [$\sim$150 cr];
(2) \textbf{twirling natywny} zamiast logicznego; (3) pełna drabina CZ, $M{=}10$--$20$ randomizacji,
$n_{\text{train}}{=}10$, wyższe shoty, $5$--$8$ punktów ewaluacji [$\sim$600--900 cr];
(4) \textbf{pasma driftu} ($\ge 2$ okna kalibracji) na~kluczowym punkcie [$\sim$200 cr];
(5) \textbf{test H3}: czy zysk twirlingu zgadza się z~przewidywaniem z~$u$. Łącznie
$\sim$1000--1300 kredytów (pełny reżim literaturowy $M{=}20$/shots$=$8192 na~całej drabinie
przekroczyłby budżet, więc skalujemy). Wynik: sprzętowa, ilościowo przewidywalna charakteryzacja
tego, kiedy twirling odzyskuje skuteczność mitygacji pod narzutem routingu.

% =====================================================================
\begin{thebibliography}{99}
% =====================================================================

% --- Mitygacja błędów: przegląd, ZNE, CDR ---
\bibitem{cai2023}
Z. Cai, R. Babbush, S. C. Benjamin, S. Endo, W. J. Huggins, Y. Li, J. R. McClean,
T. E. O'Brien, \emph{Quantum Error Mitigation}, arXiv:2210.00921 (2023).

% --- Circuit twirling / randomized compiling ---
\bibitem{wallman2016}
J. J. Wallman, J. Emerson, \emph{Noise tailoring for scalable quantum computation
via randomized compiling}, Phys. Rev. A \textbf{94}, 052325 (2016); arXiv:1512.01098.

\bibitem{hashim2021}
A. Hashim, R. K. Naik, A. Morvan, J.-L. Ville, B. Mitchell, J. M. Kreikebaum, M. Davis,
E. Smith, C. Iancu, K. P. O'Brien, I. Hincks, J. J. Wallman, J. Emerson, I. Siddiqi,
\emph{Randomized Compiling for Scalable Quantum Computing on a Noisy Superconducting
Quantum Processor}, Phys. Rev. X \textbf{11}, 041039 (2021); arXiv:2010.00215.

\bibitem{emerson2007}
J. Emerson, M. Silva, O. Moussa, C. Ryan, M. Laforest, J. Baugh, D. G. Cory,
R. Laflamme, \emph{Symmetrized characterization of noisy quantum processes},
Science \textbf{317}, 1893 (2007); arXiv:0707.0685.

\bibitem{nielsen2010}
M. A. Nielsen, I. L. Chuang, \emph{Quantum Computation and Quantum Information},
Cambridge University Press (2010) --- reprezentacja $\chi$ (macierz procesu).

\bibitem{temme2017}
K. Temme, S. Bravyi, J. M. Gambetta, \emph{Error mitigation for short-depth quantum
circuits}, Phys. Rev. Lett. \textbf{119}, 180509 (2017); arXiv:1612.02058.

\bibitem{giurgica2020}
T. Giurgica-Tiron, Y. Hindy, R. LaRose, A. Mari, W. J. Zeng, \emph{Digital zero-noise
extrapolation for quantum error mitigation}, IEEE QCE (2020); arXiv:2005.10921.

\bibitem{kandala2019}
A. Kandala, K. Temme, A. D. Córcoles, A. Mezzacapo, J. M. Chow, J. M. Gambetta,
\emph{Error mitigation extends the computational reach of a noisy quantum processor},
Nature \textbf{567}, 491 (2019); arXiv:1805.04492.

\bibitem{czarnik2021}
P. Czarnik, A. Arrasmith, P. J. Coles, L. Cincio, \emph{Error mitigation with Clifford
quantum-circuit data}, Quantum \textbf{5}, 592 (2021); arXiv:2005.10189.

% --- Readout error mitigation ---
\bibitem{maciejewski2020}
F. B. Maciejewski, Z. Zimborás, M. Oszmaniec, \emph{Mitigation of readout noise in
near-term quantum devices by classical post-processing based on detector tomography},
Quantum \textbf{4}, 257 (2020); arXiv:1907.08518.

\bibitem{nation2021}
P. D. Nation, H. Kang, N. Sundaresan, J. M. Gambetta, \emph{Scalable mitigation of
measurement errors on quantum computers}, PRX Quantum \textbf{2}, 040326 (2021);
arXiv:2108.12518.

% --- Topologia, routing, jałowe płaskowyże ---
\bibitem{li2019}
G. Li, Y. Ding, Y. Xie, \emph{Tackling the qubit mapping problem for NISQ-era quantum
devices (SABRE)}, ASPLOS (2019); arXiv:1809.02573.

\bibitem{wang2021}
S. Wang, E. Fontana, M. Cerezo, K. Sharma, A. Sone, L. Cincio, P. J. Coles,
\emph{Noise-induced barren plateaus in variational quantum algorithms}, Nat. Commun.
\textbf{12}, 6961 (2021); arXiv:2007.14384.

\bibitem{mcclean2018}
J. R. McClean, S. Boixo, V. N. Smelyanskiy, R. Babbush, H. Neven, \emph{Barren plateaus
in quantum neural network training landscapes}, Nat. Commun. \textbf{9}, 4812 (2018);
arXiv:1803.11173.

% --- Detekcja/korekcja błędów (warstwa logical-ish) ---
\bibitem{vuillot2018}
C. Vuillot, \emph{Is error detection helpful on IBM 5Q chips?}, Quantum Inf. Comput.
\textbf{18}, 949 (2018); arXiv:1705.08957.

\bibitem{linke2017}
N. M. Linke, M. Gutierrez, K. A. Landsman, C. Figgatt, S. Debnath, K. R. Brown,
C. Monroe, \emph{Fault-tolerant quantum error detection}, Sci. Adv. \textbf{3},
e1701074 (2017); arXiv:1611.06946.

\bibitem{googleqec2025}
Google Quantum AI, \emph{Quantum error correction below the surface code threshold},
Nature \textbf{638}, 920 (2025); arXiv:2408.13687.

% --- QML: benchmark, kodowanie, ansatz ---
\bibitem{bowles2024}
J. Bowles, S. Ahmed, M. Schuld, \emph{Better than classical? The subtle art of
benchmarking quantum machine learning models}, arXiv:2403.07059 (2024).

\bibitem{schuld2021}
M. Schuld, R. Sweke, J. J. Meyer, \emph{The effect of data encoding on the expressive
power of variational quantum machine learning models}, Phys. Rev. A \textbf{103},
032430 (2021); arXiv:2008.08605.

\bibitem{perezsalinas2020}
A. Pérez-Salinas, A. Cervera-Lierta, E. Gil-Fuster, J. I. Latorre, \emph{Data
re-uploading for a universal quantum classifier}, Quantum \textbf{4}, 226 (2020);
arXiv:1907.02085.

\end{thebibliography}

\end{document}
